\documentclass[../../main.tex]{subfiles}
\begin{document}

\begin{flushright}
  \footnotesize\textit{【自動文字起こし・要確認】}
\end{flushright}

% Source: handwritten.pdf 第6問．
{\bf [解]}
\begin{align}f(0)=0,\qquad f'(0)=1\label{eq:1}
\end{align}
である．任意の$a,b\in\mathbb{R}$に対し
\begin{align}
  f(a+b)=\dfrac{f(a)+f(b)}{1+f(a)f(b)},\label{eq:2}\\
  1+f(a)f(b)\ne0.\label{eq:3}
\end{align}
\cref{eq:2}で$b=-a$とすると，\cref{eq:1}から
\begin{align}f(a)+f(-a)=0.
\end{align}
だから\cref{eq:3}で$b=-a$として
\begin{align}1+f(a)f(-a)=1-\{f(a)\}^2\ne0,
\end{align}
\begin{align}f(a)\ne\pm1.
\end{align}
これが任意の$a\in\mathbb{R}$で成り立つので，$f$が連続かつ$f(0)=0$から
$-1<f(x)<1$である．

さて，\cref{eq:2}から
\begin{align}
  \begin{split}
    \dfrac{f(x+h)-f(x)}{h}
    &=\dfrac{1}{h}\left\{\dfrac{f(x)+f(h)}{1+f(x)f(h)}-f(x)\right\}\\
    &=\dfrac{f(h)}{h}\dfrac{1-\{f(x)\}^2}{1+f(x)f(h)}\\
    &\longrightarrow\{1-\{f(x)\}^2\}f'(0)=f'(x)\quad(h\to0).
  \end{split}
\end{align}
\cref{eq:1}から
\begin{align}f'(x)=1-\{f(x)\}^2.\label{eq:4}
\end{align}
したがって$y=f(x)$として，\cref{eq:4}をあわせて微分方程式
\begin{align}\dfrac{1}{1-y^2}\,dy=dx
\end{align}
が成り立つ．これをといて，$-1<y<1$から
\begin{align}\log\dfrac{1+y}{1-y}=2x+C.
\end{align}
$x=0$で$y=0$であることから，$C=1$であり，
% 原稿には「+C」「C=1」とある．ここは原稿内で不整合だが，後続式に合わせて黙って訂正しない．
\begin{align}
  \dfrac{1+y}{1-y}=e^{2x},\qquad y=\dfrac{e^{2x}-1}{e^{2x}+1}.
\end{align}
逆に，このとき\cref{eq:1,eq:3}はみたし，\cref{eq:2}については
\begin{align}
  \begin{split}
    \dfrac{f(a)+f(b)}{1+f(a)f(b)}
    &=\dfrac{(e^{2a}+1)(e^{2b}-1)+(e^{2a}-1)(e^{2b}+1)}{(e^{2a}+1)(e^{2b}+1)+(e^{2a}-1)(e^{2b}-1)}\\
    &=\dfrac{2(e^{2a+2b}-1)}{2(e^{2a+2b}+1)}=f(a+b)
  \end{split}
\end{align}
から成立し十分．以上から
\begin{align}f(x)=\dfrac{e^{2x}-1}{e^{2x}+1}
\end{align}
であり，
\begin{align}f''(x)=\{1-\{f(x)\}^2\}'=-2f(x)f'(x).
\end{align}
これは$x>0$で$f(x)>0$，$f'(x)>0$から$f''(x)<0$となり，したがって上に凸である．

\end{document}
